% LaTeX fragment. Requires amsmath, amssymb, amsthm, hyperref.
% This file distinguishes literal text, reconstruction, and this run's proofs.
% No statement from OBJ_H035_results.zip was used to produce this draft.

\section*{Conventions and proof dependencies}
All rings are commutative with identity; the zero ring is admitted. A finite
module means finitely generated. A strict inverse system has surjective
transition maps; ``stationary'' means that all transitions beyond a stage
are isomorphisms, not merely that the pro-object is isomorphic to a constant
one. A representable covariant module functor means
$M\mapsto\operatorname{Hom}_A(P,M)$. All reconstructions below are ours.
The manuscript's numbered lemmas are not silently treated as established.
In particular H035.10, H035.15 and the literal H035.29 cannot be used as
written. H035.30 uses the corrected, explicitly proved versions.

\section{H035.1: the title-sheet programme}
\textbf{His text.} The title proposes prorepresentation of a functor
$G/S$ along a section and invariants denoted by an $\Omega^1$-expression.
It specifies neither a category of test objects nor representability
hypotheses. It is a programme, not an assertion that every functor works.

\textbf{Our precise construction.} For an augmented $A$-algebra
$C\xrightarrow{e}A$, put $I=\ker(e)$. On augmented $A$-algebras $T$ with
nilpotent augmentation ideal, the maps $C\to T$ which send $I$ into that
ideal form the filtered union of the functors represented by
$C/I^{n+1}$. Indeed, such a map kills $I^{n+1}$ once the target ideal has
nilpotence exponent at most $n+1$, and the converse is immediate. Thus
the formal completion along the section is prorepresented by this system.
For $T=A\oplus M$, with $M^2=0$, a map lifting $e$ is uniquely $e+d$,
where $d$ is an $A$-derivation. Its restriction to $I$ kills $I^2$;
conversely an $A$-linear map $I/I^2\to M$ extends uniquely to such a
derivation, because $C=A\oplus I$ as an $A$-module. Hence the tangent
functor is $\operatorname{Hom}_A(I/I^2,M)$, and
$I/I^2\simeq\Omega^1_{C/A}\otimes_C A$. Multiplication gives a
surjection $\operatorname{Sym}_A(I/I^2)\twoheadrightarrow
\bigoplus I^n/I^{n+1}$, since products of elements of $I$ generate $I^n$.
These are constructions attached to the stated augmented algebra; they
are not additional words supplied to the title sheet.

\textbf{Why an unrestricted extension is false.} Over a field of
characteristic different from two, take the functor sending an augmented
algebra $T$ to the elements $z$ in its augmentation ideal with $z^2=0$,
modulo $z\sim-z$. It has the zero section. On square-zero algebras its
tangent functor is $M\mapsto M/\{\pm1\}$. The distinct classes of
$(1,1)$ and $(1,-1)$ in $(k\oplus k)/\{\pm1\}$ have the same image
in $(k/\{\pm1\})^2$. Thus it does not preserve this fibre product.
Every represented functor preserves limits, and a filtered colimit of
represented functors preserves finite limits in sets. This functor is
therefore not prorepresentable. A bare title cannot authorize a universal
prorepresentability theorem.

\section{H035.2: restriction of scalars and the linear structure}
\textbf{Our proof.} Restricting a represented functor to $B$-modules gives
$\operatorname{Hom}_A(P,N)=\operatorname{Hom}_B(B\otimes_A P,N)$ by
tensor--Hom adjunction. For a finite-limit-preserving set-valued functor
$F$ on modules, $F(0)$ is a singleton and
$F(M\oplus M)=F(M)\times F(M)$. Apply $F$ to addition, zero, negatives
and multiplication by each $a\in A$. All module identities are diagrams
of module maps and finite products, so they remain identities after $F$.
Thus $F(M)$ has a canonical $A$-module structure, functorially, and $F$
is $A$-linear. In particular it preserves kernels and is left exact as a
module-valued functor. This explains all subsequent tensor comparisons;
mere values in sets would not explain them without these hypotheses.
If $F$ commutes with filtered colimits, the expression of a localization
as a filtered colimit of multiplication maps gives
$F(M)_p\simeq F(M_p)$. More generally, writing a flat module as a
filtered colimit of finite free modules gives
$F(M)\otimes_A B\simeq F(M\otimes_A B)$ for a flat $A$-algebra $B$.
The localization assertion itself needs only the explicit multiplication
colimit, not the general flat-module theorem.

\section{H035.3: faithfully flat representability descent (Lemme 1)}
\textbf{Our proof.} Necessity follows from H035.2 and left exactness of Hom
in its second argument. Conversely suppose $F$ is left exact and
$F_B=\operatorname{Hom}_B(Q,-)$ for faithfully flat $A\to B$. Restrict
along the two maps $B\rightrightarrows B\otimes_A B$. Both resulting
functors are the same restriction of $F$. Yoneda gives an isomorphism
between the two pullbacks of $Q$, with the cocycle identity over
$B\otimes_A B\otimes_A B$, because the identifications of functors
are identities. Effective faithfully flat descent for modules gives
$Q\simeq B\otimes_A P$ with this datum.

For completeness, the descended module is the equalizer of the two
semilinear maps on $Q$ prescribed by the datum. Tensor its defining
equalizer with $B$; the datum and the multiplication map on $B$ supply
the usual contracting map, identifying the tensor equalizer with $Q$.
Faithful flatness then verifies the original equalizer and its uniqueness.
For any $M$, faithfully flat descent expresses $M$ as the equalizer of
$B\otimes M\rightrightarrows B\otimes B\otimes M$. Applying $F$ and
using left exactness identifies $F(M)$ with the equalizer which computes
$\operatorname{Hom}_A(P,M)$. The construction is natural in $M$.

Finite generation descends: finitely many generators of $B\otimes P$
involve only finitely many elements of $P$; if they generate from a finite
submodule $P_0$, then $B\otimes(P/P_0)=0$, whence $P=P_0$. For finite
presentation choose a finite free surjection onto $P$. Flatness identifies
the base-changed kernel, which is finite if $Q$ is finitely presented;
the just-proved finite-generation descent makes the original kernel
finite. This proves all three versions, not only the noetherian one.

\section{H035.4: gluing a principal completion and its complement (Lemme 2)}
\textbf{Precise reconstruction.} Let $A$ be noetherian, $f\in A$,
$A_n=A/(f^n)$ for $n\geq1$, and $\widehat A=\lim A_n$. Assume left
exactness, preservation of filtered colimits, the comparison
$F(M)\simeq\lim F(M/f^nM)$ for every finite $\widehat A$-module $M$,
and representability of $F_{A_f}$ and every $F_{A_n}$. Then $F$ is
represented by a finite $A$-module; the converse holds.

\textbf{Our proof.} A module representing a functor commuting with
filtered colimits is finitely presented: write the module as a filtered
colimit of finitely presented modules and factor its identity through a
stage. It is a retract of that finitely presented stage, hence finitely
presented. Over a noetherian ring this is equivalent to finiteness.
Let $P_n$ represent $F_{A_n}$. Yoneda gives compatible isomorphisms
$P_{n+1}/f^nP_{n+1}\simeq P_n$ and surjective transition maps.

Here are the details of the completion step. Choose generators of $P_1$
and lift them successively to compatible elements in $P_n$. Nakayama over
$A_n$, where $(f)$ is nilpotent, says that these $r$ elements generate
every $P_n$. The kernels of $A_n^r\to P_n$ have surjective transition
maps: lift a kernel element, and correct its image in
$\ker(P_{n+1}\to P_n)=f^nP_{n+1}$ by a vector in $f^nA_{n+1}^r$.
Taking inverse limits gives a surjection $\widehat A^r\to P=\lim P_n$.
Its kernel is closed and finite, since $\widehat A$ is noetherian.
Write this kernel as $K$. Surjectivity of the kernel transitions says
that $K$ maps onto the kernel of $A_n^r\to P_n$. Consequently
$\ker(\widehat A^r\to P_n)=K+f^n\widehat A^r$, and hence
$P_n=\widehat A^r/(K+f^n\widehat A^r)=P/f^nP$.
Thus $P$ is finite and complete. This does not assert that the original
kernel system is itself the system $K/f^nK$.

For finite complete $M$, the assumed comparison and representability give
\[
 F(M)=\lim_n\operatorname{Hom}_{A_n}(P_n,M/f^nM)
       =\operatorname{Hom}_{\widehat A}(P,M).
\]
The last equality follows by compatible reduction of maps, with existence
and uniqueness from completeness. Every $\widehat A$-module is the
filtered union of its finite submodules, so the equality holds for all
modules. The algebra $\widehat A\times A_f$ is faithfully flat over
$A$: completion is flat; primes containing $f$ lift to the completion
and the other primes are covered by $A_f$. H035.3 now descends the finite
representing module. Conversely Hom commutes with inverse limits and
with these base changes, and a finite representing module commutes with
filtered colimits. This proves necessity including the exact completed
test modules written in the scan.

\section{H035.5: the noetherian-induction criterion (Proposition 3)}
\textbf{Reconstruction.} Quotients in the generic nonnilpotent-element
condition must be nonzero. If the zero quotient is included, the literal
necessity assertion is false even for $F=0$, since the zero ring has no
nonnilpotent element. The corrected condition is: every nonzero quotient
$B$ has a nonnilpotent $g$ with $F_{B_g}$ representable. Together with
left exactness, filtered colimits, and every principal-completion
comparison, it is equivalent to finite representability.

\textbf{Our proof.} The conditions pass to quotient rings. For completion
this uses $\widehat{A/I}=\widehat A/I\widehat A$ and the fact that a
finite module over the former is finite over the latter. If the result
fails choose, by the ascending chain condition, an ideal $I$ maximal
among those for which $F_{A/I}$ is not represented by a finite module.
Replace $A$ by $A/I$. Choose the stipulated nonnilpotent $f$. Every
$(f^n)$ is nonzero, so maximality gives representability over every
$A/(f^n)$. H035.4 supplies representability over $A$, a contradiction.
Necessity follows by restricting a finite representing module; any
nonzero quotient permits $f=1$. This proof also covers the zero ring
separately, where the zero functor is the only left-exact functor.

\section{H035.6: the complete-local criterion (Proposition 4)}
\textbf{Our proof.} A noetherian local ring has faithfully flat maximal-
ideal completion. The given completion comparisons reduce the assertion
to a complete local ring. Put $A_n=A/\mathfrak m^n$. Let $E_n$ be the
injective hull of its residue field. Artinian Matlis duality says that
$E_n$ is finite, $\operatorname{End}_{A_n}(E_n)=A_n$, and
$D_n=\operatorname{Hom}_{A_n}(-,E_n)$ is an exact duality on finite
modules. These are the artinian case of Stacks, Proposition 47.7.8
(the exact source is recorded in the reference audit).

The finite-value hypothesis makes $F(E_n)$ finite. Set
$P_n=D_n(F(E_n))$. Every finite $A_n$-module has an injective
copresentation $0\to M\to E_n^r\to E_n^s$: embed $M$ in a finite
sum of copies of the injective cogenerator, then embed the finite
cokernel. On such sums, the canonical identification
$F(E_n^r)=\operatorname{Hom}_{A_n}(P_n,E_n^r)$ is compatible with
all matrices, since $\operatorname{End}(E_n)=A_n$ and $F$ is linear.
Taking the kernel of the copresentation gives the identification for
$M$, naturally; filtered colimits give it for every module. The $P_n$
are compatible by Yoneda. Apply the complete-module gluing argument of
H035.4 with $\mathfrak m$ in place of $(f)$; its proof uses powers of a
finitely generated ideal and is unchanged. The assumed inverse-limit
comparison identifies $F$ with the resulting finite Hom functor.
Faithfully flat descent returns to the original ring.

The parenthetical weakening in the manuscript is valid: the proof needs
finite values only on the finite $A$-modules $E_n$, which are already
complete. Necessity is immediate from finite Hom and completeness.

\section{H035.7: the local representations used in Lemme 5}
\textbf{Our proof.} Assume Lemme 5(i)--(iv). At a prime $p$, preservation
of filtered colimits gives the localization identity of H035.2.
The stated comparison is exactly the completion condition of H035.6
for $A_p$. Its artinian injective hulls have finite length over $A_p$,
so (iv), which asks only for finite-length values, suffices. Consequently
$F_{A_p}=\operatorname{Hom}_{A_p}(Q_p,-)$ with $Q_p$ finite. If
$q\subset p$, restricting this identification to $A_q$-modules gives
$Q_p\otimes_{A_p}A_q\simeq Q_q$, uniquely. No generic-representation
condition is needed for this local assertion.

\section{H035.8: strict prorepresentation and minimal supports (Lemme 5)}
\textbf{Reconstruction.} In (v) require $0\ne f\in A/p$. The enlarged
scan does not display a legible nonzero qualification. Allowing $f=0$
makes (v) vacuous and the counterexample in H035.10 refutes the conclusion.
We now prove the nonzero version in full.

\textbf{Our proof of the support assertion.} For a finite $M$ and
$\alpha\in F(M)$, define
$N(p)=\operatorname{im}(Q_p\xrightarrow{\alpha_p}M_p)$.
These submodules commute with generization. Consider opens $U$ on which
they are the stalks of a coherent submodule $N_U\subset\widetilde M$.
They agree on overlaps, so their union has this property; it is coherent
locally. There is therefore a largest such $U$, possibly initially empty.
Suppose $U\ne\operatorname{Spec}A$, and choose a generic point $p$ of
its closed complement. Shrink an affine neighbourhood $V$ of $p$ so that
the complement in $V$ lies in $V(p)$.

There is a coherent submodule $N'\subset\widetilde M|_V$ whose stalk
at $p$ is $N(p)$ and whose restriction to $U\cap V$ is $N_U$.
Here is the spreading justification. Extend $N_U$ to a finite submodule
locally in $\widetilde M$, spread generators of $N(p)$, and compare their
sum and intersection. On the punctured $\operatorname{Spec}A_p$ the
two submodules agree, by generization. The finitely many irreducible
components of the supports of their difference which meet $U$ do not
contain $p$: a generic point of such a component lying below $p$ would
belong to the punctured spectrum where equality holds. Removing those
components gives the required agreement near $p$.
The image of $\alpha$ in $F(M/N')$ vanishes at $p$. Localization of
$F$ allows another shrinking for which it vanishes throughout $V$.
Left exactness then gives $\alpha\in F(N')$. Replace $M$ by $N'$.
Now $N(p)=M_p$, and equality already holds on $U\cap V$.

Condition (v) represents the functor over a nonempty open of
$\operatorname{Spec}(A/p)$ by a finite module. The image of
$\alpha$ modulo $p$ corresponds to a map from that representing
module to $M/pM$. At the generic point it is surjective because
$N(p)=M_p$. The cokernel is finite, so it remains surjective on a
smaller neighbourhood of that generic point. At every $s\in V(p)$
in this neighbourhood, its image is
$(N(s)+pM_s)/pM_s$, by restriction and Yoneda. Nakayama gives
$N(s)=M_s$. Thus all supports are coherent on a larger open than $U$,
a contradiction. They therefore give a global finite submodule $N$.
The image of $\alpha$ in $F(M/N)$ vanishes at every prime, hence is
zero; left exactness shows $\alpha\in F(N)$. Every submodule defining
$\alpha$ contains $N$ stalkwise, so $N$ is the unique minimal support.

\textbf{Our construction of the pro-object.} Use the finite pairs
$(P,\xi)$ with minimal support equal to $P$. A map of pairs has
surjective underlying map: its image is a defining submodule for the
target element. Two parallel maps are equal, because the source element
belongs to the kernel of their difference by left exactness, and has
full support. Finite products and equalizers, followed by taking the
minimal support, show that this category is cofiltered; it is essentially
small. Every element of $F(M)$ comes from a finite submodule of $M$
and then from its minimal support. If two such presentations give the
same element, the fibre product over $M$, with its induced element,
and then its minimal support identify them. Therefore
\[
 F(-)=\mathop{\rm colim}_i\operatorname{Hom}_A(P_i,-)
\]
with all transition maps $P_j\twoheadrightarrow P_i$ surjective.

At $p$ this union of subfunctors is represented by the finite $Q_p$ of
H035.7. The identity of $Q_p$ occurs at some stage. By Yoneda the
inclusion of that stage is induced by an epimorphism $Q_p\to(P_i)_p$;
factoring the identity supplies a left inverse, so it is an isomorphism.
The same reasoning applies to every later stage and proves eventual
stationarity of the localized system, as claimed.

\section{H035.9: generic representation implies constructible fibre rank}
\textbf{Our proof.} On every integral closed subset $\operatorname{Spec}B$
condition (v) gives a nonempty open where $F$ is represented by a finite
module $P$. On a smaller nonempty open $P$ is free: choose a basis over
the fraction field, spread it to a map from a finite free module, and
kill the finite torsion kernel and cokernel by one nonzero element.
Thus $\dim_{k(p)}F(k(p))$ is constant there. Noetherian induction on
the complementary closed subsets produces a finite stratification by
locally closed subsets on which the rank is constant. This is precisely
constructibility of the integer-valued function. The generically normal
assumption written in the insertion is unnecessary for this direction.

\section{H035.10: the false converse of Remarque 5 bis and its repair}
\textbf{Our counterexample.} Put $A=\mathbb Z$ and
$L=\sum_{p\ {
m prime}}p^{-1}\mathbb Z\subset\mathbb Q$.
The elements of $L$ are exactly the rationals with squarefree
denominator: Bezout expresses $1/(p_1\cdots p_r)$ as an integer
linear combination of the $1/p_i$, and sums introduce no squared
denominator. Let $F(M)=L\otimes M$. Torsion-free modules over $\mathbb Z$
are flat, so $F$ is left exact, and tensor commutes with filtered colimits.
At $(p)$, $L_{(p)}=p^{-1}\mathbb Z_{(p)}$, a free rank-one module;
at $(0)$ it is $\mathbb Q$. Hence all the local functors are finite
rank-one representables. For finite $\widehat{\mathbb Z_{(p)}}$-modules,
their completion comparisons are those of the identity functor, and
finite-length values are finite. The corresponding assertions at $(0)$
are immediate. Every residue-field rank is one, a constructible function.
All integral quotients of $\mathbb Z$ are normal.

For any nonzero integer $f$, $L_f$ is not finite over $\mathbb Z_f$:
a proposed finite set of generators uses finitely many denominator
primes; choose $q$ outside that list and outside the prime divisors of
$f$. Its span cannot contain $1/q$. If $F_f$ were represented, filtered-
colimit preservation would make its representing module finitely
presented, and exactness would make that module projective. Evaluation
at $\mathbb Z_f$ would then make $L_f$ finite, a contradiction. Thus
(v bis) does not imply (v), with every displayed hypothesis checked.
The very same example disproves Lemme 5 if its missing nonzero
qualification is read literally: strict prorepresentation would give
a nonzero finite subfunctor of $L\otimes-$, which is impossible by
the minimal-support argument, or directly the element $1\in L=F(A)$
has defining ideals $pA$ for all primes $p$, whose intersection is zero.
Indeed $1=p(1/p)$ lies in the image $L\otimes pA\to L$, whereas
$F(0)=0$. No minimal defining submodule exists.

\textbf{Our corrected theorem.} In addition to Lemme 5(i)--(iv), suppose
$F(B)$ is finite for every integral quotient $B$ of $A$. Then the
constructible rank condition implies (v), with no normality hypothesis.
Choose an integral $B$ and a nonempty open on which the rank is a
constant $n$. For every point there, let $P_p$ be its finite local
representing module. It has generic rank $n$, and
$\dim_{k(p)}P_p/pP_p=n$ by the fibre-rank condition. Nakayama gives
a surjection $B_p^n\to P_p$; its kernel has rank zero and is a
submodule of a free module over the domain $B_p$, hence is zero.
Thus $P_p$ is free. The finite module $V=F(B)$ localizes to
$\operatorname{Hom}(P_p,B_p)$, so it is locally free of rank $n$
on this open. The canonical evaluation map
$V\otimes_B M\to F(M)$ is an isomorphism at every point, because
locally the functor is represented by a finite free module. It is
therefore an isomorphism for every $M$. The representing module is
$V^\vee$. This proves the repaired implication required by Théorème 12,
whose condition (iv) supplies the extra hypothesis. The manuscript's
use of ``$F(N)$ est de type fini'' is exactly the missing step.

\section{H035.11: generic freeness of the nilpotent graded modules (Lemme 6a)}
\textbf{Our proof.} Work locally on a noetherian scheme. Its nilradical
$J$ is nilpotent and its associated graded algebra and the associated
graded of a coherent module have finitely many coherent pieces over
$A/J$. Remove intersections of irreducible components of the reduced
scheme, then use generic freeness on each integral component for this
finite list of modules. The resulting open is dense; the graded pieces
are locally free there. On an integral affine component generic freeness
for a finite module is proved by choosing a fraction-field basis and
killing the finite torsion kernel and cokernel. The locally noetherian
version follows by taking the union of the local good opens, since the
property is local. This proves both the algebra and module assertions.

\section{H035.12: regular sequences and all positive Tor (Lemme 6b)}
\textbf{Our filtered proof.} At a point of the open in H035.11 write
$A_0=A/J$. Every graded piece of $A$ and of the module $R$ is free over
the local ring $A_0$. If $x$ is $A$-regular, its image $\bar x$ is
$A_0$-regular. Otherwise a nonzero annihilator of $\bar x$, multiplied
by a basis vector of the last nonzero $J^r/J^{r+1}=J^r$, would give
a nonzero element annihilated by $x$ in $A$. Consequently $x$ acts
injectively on every graded piece of $R$, so on $R$ itself. Moreover
$xR\cap J^aR=xJ^aR$: if $xr\in J^aR$ and $r$ has a smaller
filtration degree, its nonzero initial form is killed by $\bar x$,
a contradiction. Thus
\[
 \operatorname{gr}_J(R/xR)=\operatorname{gr}_J(R)/\bar x,
 \qquad
 \operatorname{gr}_J(A/xA)=\operatorname{gr}_J(A)/\bar x.
\]
These are free graded modules over $A_0/(\bar x)$. The image of $J$
is still nilpotent, even if it is not the whole nilradical of the quotient.
Repeat with an $A$-regular sequence $x_1,\ldots,x_r$; it is $R$-regular.
Its Koszul complex is a finite free resolution of $A/(\mathbf x)$,
and its tensor with $R$ is exact in positive homological degrees.
Hence $\operatorname{Tor}_i^A(R,A/(\mathbf x))=0$ for every $i>0$.
The undefined $M$ in the manuscript proof is not adopted as a new object;
our $R$ is the coherent module to which the lemma applies.

\section{H035.13: generic Hom base change (Corollaire 7)}
\textbf{Our proof.} On an affine noetherian open choose a finite free
presentation $L_1\to L_0\to P\to0$ and form
$K=[L_0^\vee\otimes Q\to L_1^\vee\otimes Q]$ in degrees zero and
one. Then $H^0K=\operatorname{Hom}(P,Q)$, and for every $M$,
$H^0(K\otimes M)=\operatorname{Hom}(P,Q\otimes M)$. Put
$Z=\ker d$, $B=\operatorname{im}d$, $H=\operatorname{coker}d$.
Apply H035.12 to the finite list $B,H$ (or to $K^1,H$ with all positive
Tor). After a common dense shrinking,
$\operatorname{Tor}_1(B,M)=\operatorname{Tor}_1(H,M)=0$ for all
the regular-sequence quotients in question. Tensor the two sequences
$0\to Z\to K^0\to B\to0$ and
$0\to B\to K^1\to H\to0$. They show that the kernel of
$K^0\otimes M\to K^1\otimes M$ is exactly $Z\otimes M$.
This is the required comparison. The construction is local, so these
dense good opens give the asserted sheaf isomorphism. It repairs the
written proof without weakening the conclusion.

\section{H035.14: the coefficient in the displayed Hom complex}
\textbf{Literal computation and counterexample.} The scan has $P$ in
both tensor factors. With that coefficient the displayed complex computes
$\operatorname{Hom}(P,P)$ in degree zero. For example, over a field,
$P=k$, $Q=k^2$, $L_0=k$, $L_1=0$, the displayed $P$-complex has
one-dimensional $H^0$, whereas $\operatorname{Hom}(P,Q)$ has dimension
two. Thus it cannot be the stated computation.

\textbf{Corrected computation.} Apply $\operatorname{Hom}(-,Q)$ to
$L_1\to L_0\to P\to0$. Left exactness gives
$\operatorname{Hom}(P,Q)=\ker(\operatorname{Hom}(L_0,Q)\to
\operatorname{Hom}(L_1,Q))$. Since $L_i$ are finite free,
$\operatorname{Hom}(L_i,Q)=L_i^\vee\otimes Q$. This proves exactly
the correction used in H035.13. The ink is retained in READINGS.tsv.

\section{H035.15: the two-Tor criterion is false as written}
\textbf{Our counterexample.} Let $A=k[x,y]/(xy)$, $M=A/(x)$ and
$K=[A\xrightarrow{y}A]$ in degrees zero and one. Then
$H^0K=(x)$, $H^1K=A/(y)$ and $Z^1K=A$. A free resolution of $M$
has successive differentials $x,y,x,y,\ldots$. It is exact because
$\operatorname{ann}(x)=(y)$ and $\operatorname{ann}(y)=(x)$ in $A$.
Tensoring with $A/(y)=k[x]$, the differential into degree zero is the
injective multiplication by $x$, so $\operatorname{Tor}_1(M,H^1K)=0$.
Also $\operatorname{Tor}_1(M,Z^1K)=0$ since $Z^1K=A$.
But
\[
 H^0K\otimes M=(x)/(x^2)\simeq k,\qquad
 H^0(K\otimes M)=\ker(y:k[y]\to k[y])=0.
\]
Thus the comparison is not an isomorphism despite both written
vanishings. This counterexample concerns arbitrary $M$, as in the
criterion; its $M$ is not a quotient by an $A$-regular sequence.

\textbf{Corrected criterion and proof.} With
$B^1=\operatorname{im}(K^0\to K^1)$, the sufficient conditions are
$\operatorname{Tor}_1(M,B^1)=\operatorname{Tor}_1(M,H^1K)=0$,
by the two exact sequences used in H035.13. A condition avoiding $B^1$
is
\[
\operatorname{Tor}_1(M,K^1)=\operatorname{Tor}_1(M,H^1K)
=\operatorname{Tor}_2(M,H^1K)=0.
\]
Indeed the long Tor sequence of $0\to B^1\to K^1\to H^1K\to0$
then gives $\operatorname{Tor}_1(M,B^1)=0$. H035.12 supplies all
these vanishings in the actual application. The subsequent argument is
therefore valid after repair, even though its general sufficient
criterion is false.

\section{H035.16: the heading and the canonical-module criterion}
\textbf{His words.} ``module de différentielles pour un C.M.''

\textbf{Reading and reconstruction.} The enlarged scan says
\emph{différentielles}; it does not say \emph{dualité}. No base ring for
Kähler differentials is specified. We do not attribute a modern Kähler
definition to this heading. The property immediately written after it,
and the subsequent \emph{faisceau fondamental}, motivate the following
explicitly modern reconstruction. If $(A,\mathfrak m,k)$ is a noetherian
Cohen--Macaulay local ring and $W$ is finite, then $W$ is a canonical
module if and only if for every system of parameters $\mathbf x$,
$W/\mathbf xW$ is the injective hull of $k$ over $A/(\mathbf x)$.
Canonical here means that its completion is the canonical module of
$\widehat A$. If one instead reads the heading as a claim about ordinary
Kähler differentials, it is false: $\Omega^1_{k/k}=0$, whereas the
injective hull required for the zero-dimensional local ring $k$ is $k$.

\textbf{Our proof of the reconstruction.} The forward implication is
H035.17--18. Conversely fix a system of parameters $x_1,\ldots,x_d$.
For every choice of positive integers $n_i$, artinian Matlis duality gives
\[
 \ell\bigl(W/(x_1^{n_1},\ldots,x_d^{n_d})W\bigr)
 =\ell\bigl(A/(x_1^{n_1},\ldots,x_d^{n_d})\bigr)
 =n_1\cdots n_d\ell\bigl(A/(\mathbf x)\bigr).
\]
The last equality follows successively from regularity of the parameters
on the Cohen--Macaulay ring $A$. Fix $n_2,\ldots,n_d$ and put
$N=W/(x_2^{n_2},\ldots,x_d^{n_d})W$. The displayed length formula is
$\ell(N/x_1^nN)=nC$. Multiplication by $x_1$ from $N/x_1^nN$ to
$N/x_1^{n+1}N$ has cokernel $N/x_1N$, of length $C$, so has image
of length $nC$ and is injective. Hence $\operatorname{ann}_N(x_1)$
lies in every $x_1^nN$ and is zero by Krull intersection. In particular
$\operatorname{ann}_W(x_1)$ lies in every
$(x_2^m,\ldots,x_d^m)W$, and is zero by the same theorem. The quotient
$W/x_1W$ satisfies the identical parameter-quotient test over $A/x_1A$.
Induction proves that $W$ is maximal Cohen--Macaulay.

Pass to the completion, where a canonical module $\omega$ exists.
Maximal Cohen--Macaulayness and the test for the chosen original
parameter system persist under this faithfully flat extension. We do
not need to assert a test for newly chosen formal parameter systems.
Canonical duality, Stacks Lemma 47.16.7 (tag 0B5A), applied to the
normalized complex $\omega[d]$, gives a duality
$D=\operatorname{Hom}(-,\omega)$ on maximal Cohen--Macaulay modules,
with positive Ext vanishing and biduality. Let $L=D(W)$. Reduction
modulo the regular system $\mathbf x$, using that Ext vanishing, gives
\[
 L/\mathbf xL\simeq
 \operatorname{Hom}_{A/(\mathbf x)}(W/\mathbf xW,
                                      \omega/\mathbf x\omega)
 \simeq\operatorname{End}(E)=A/(\mathbf x).
\]
The last equality is artinian Matlis duality, Stacks Proposition 47.7.8
(08Z9). Nakayama gives a surjection $A\to L$. Its kernel $K$ satisfies
$K/\mathbf xK=0$: tensor the exact sequence with $A/(\mathbf x)$ and
use that $\mathbf x$ is $L$-regular, so its first Tor vanishes.
Nakayama gives $K=0$, hence $L=A$ and $W=D(L)=\omega$. This proves
the stated criterion, including the descent interpretation through
completion. In dimension zero it is precisely artinian Matlis duality.

\section{H035.17: fundamental modules and regular quotients (Lemme 8)}
\textbf{Our modern statement and proof.} Let $A$ be a noetherian
Cohen--Macaulay local ring possessing a canonical module $\omega$,
and let $x_1,\ldots,x_r$ be an $A$-regular sequence in its maximal
ideal. Then it is $\omega$-regular and
$\omega/(\mathbf x)\omega$ is a canonical module for
$A/(\mathbf x)$. This statement localizes and thus gives the asserted
sheaf statement. To prove it, completion is faithfully flat and detects
regularity and canonical modules, so assume $A$ complete. Its normalized
dualizing complex is $\omega[d]$ (Stacks 0DW4). The finite-morphism
duality construction gives
$R\operatorname{Hom}_A(A/(\mathbf x),\omega[d])$ as the normalized
dualizing complex of the quotient. The canonical module is maximal
Cohen--Macaulay by 0DW4, so the regular sequence is also regular on it.
The dual Koszul complex therefore has its only cohomology in degree $r$,
equal to $\omega/(\mathbf x)\omega$. After the shift by $d$, this is
the normalized complex of a Cohen--Macaulay ring of dimension $d-r$.
It follows that this quotient is the canonical module. These local
identifications are compatible with restriction, so glue to the claimed
fundamental sheaf on the closed subscheme. The finite duality and CM
facts are exact background hypotheses, not an assertion that an arbitrary
locally noetherian scheme has a global dualizing complex.

\section{H035.18: the parameter quotient (Corollaire 9)}
\textbf{Our proof.} In a Cohen--Macaulay local ring every system of
parameters is a regular sequence. Apply H035.17 with $r=d$.
The resulting ring is artinian local. Its normalized dualizing module
is $E(k)$: its Matlis dual is the ring, and Matlis biduality identifies
every canonical module with $E(k)$. It is injective by definition of
the injective hull, indecomposable because its socle is the simple
essential submodule $k$, and a cogenerator of finite modules because
each nonzero finite module has a nonzero map to it. Conversely its
endomorphism ring is the artinian ring itself by 08Z9. These observations
justify exactly the injective and detection properties used later.

\section{H035.19: generic canonical modules (Corollaire 10)}
\textbf{Precise statement.} If $X$ is locally noetherian and locally
embeddable in a regular scheme, there are a dense open $U\subset X$
and a coherent $\mathcal O_X$-module $\Omega$ such that $U$ is
Cohen--Macaulay and $\Omega_x$ is a canonical module for every $x\in U$.
Its parameter quotients have the properties of H035.18, and the Hom
detection property is proved in H035.27. The same conclusion holds if
$X$ is Gor-2: every integral closed subscheme locally contains a nonempty
Gorenstein open. This stronger modern sufficient hypothesis is ours.

\textbf{Exact source and reduction.} Lyu, \emph{Formal lifting of
dualizing complexes and consequences} (2025), Lemma 6.1(ii), says:
if $A$ is noetherian, $p\in\operatorname{Spec}A$, $A/p$ has a
nonempty Gorenstein open, and a finite module $M$ has $M_p$ canonical
over $A_p$, then $M_q$ is canonical for all
$q\in V(p)\cap D(f)$ for some $f\notin p$. The full statement and
proof were opened. Corollary 4.7 there says Gorenstein rings are Gor-2
and their finite type algebras are Gor-2. Thus quotients of regular
rings satisfy precisely the needed hypothesis; no finite global Krull
dimension assumption on the regular ring is inserted.

For completeness, here is the construction from that exact spreading
lemma. Near a generic point $\eta$, the ring $A_\eta$ is artinian and
has a finite canonical module, its residue-field injective hull. A finite
presentation over $A_\eta$ spreads, by clearing finitely many
denominators, to a finite module on an affine neighborhood. Lyu's lemma
then makes this module canonical on an open of that component. Remove
the other components near $\eta$. The CM locus is open here by EGA IV
6.11.8, because Gor-2 implies its CMU hypothesis; it contains $\eta$.
We may therefore shrink once more so that all points are CM. Choose
one such open in each irreducible component, disjoint from all the other
components. These opens are pairwise disjoint, and their union is dense.
The finite modules on them glue, with no transition ambiguity, to a
coherent module on this open. Extend it coherently to $X$ as follows.
On an affine noetherian scheme $W$, a coherent sheaf $E$ on an open
$V$ is a restriction of a finite module: $V$ is quasi-compact, its
direct image is quasi-coherent, and finitely many elements of the
corresponding module suffice to generate $E$ on a finite principal
cover of $V$. Their finite submodule restricts to $E$.
For a locally noetherian $X$, order coherent extensions of $E$ to
larger opens, with the identifications fixed. Chains glue, so Zorn
gives a maximal extension. If its domain is not $X$, take an affine
noetherian $W$ meeting the complement. The affine result extends
the existing sheaf from the open intersection to $W$, and gluing
enlarges the domain, a contradiction. This proves the extension
without global quasi-compactness. Values outside $U$ are immaterial.

\section{H035.20: the excellent version of the margin}
\textbf{Our proof and scope.} An excellent noetherian ring is J-2;
the regular locus of every finite type algebra over it is open. In
particular an integral quotient has a nonempty regular open, since its
generic local ring is a field. A regular ring is Gorenstein. Thus an
excellent scheme is Gor-2, and H035.19 applies. This proves the
excellent clause of the margin. It asserts a coherent generic canonical
module; it does not assert that every excellent scheme possesses a
global dualizing complex. Lyu's Remark 7.9 explicitly warns against the
latter assertion.

\section{H035.21: the CMU-only extension is false}
\textbf{His words.} ``ou seulement qu'il satisfasse les conditions de
EGA IV 6.11.8.''

\textbf{The cited condition.} The opened Proposition EGA IV 6.11.8,
printed p.162, assumes that every integral closed subscheme contains a
nonempty Cohen--Macaulay open (CMU). Its conclusions concern codepth
and CM loci. It does not supply canonical modules.

\textbf{Our counterexample, with noetherianity checked.} Fix an infinite
field $k$ and algebraically independent $t_1,t_2,\ldots$. Set
\[
 B=k[t_i^3,t_i^4,t_i^5:i\geq1],\qquad
 \mathfrak m_i=(t_i^3,t_i^4,t_i^5)B,\qquad
 S=B\setminus\bigcup_i\mathfrak m_i,\qquad A=S^{-1}B.
\]
Each $\mathfrak m_i$ is prime, since its quotient is the domain in
the remaining variables, so $S$ is multiplicative. A nonzero element
of $B$ involves finitely many blocks and belongs to only finitely many
$\mathfrak m_i$.

We first show that every prime disjoint from $S$ is contained in one
$\mathfrak m_i$. Suppose otherwise and choose $0\ne x$ in it. The
set $T=\{i:x\in\mathfrak m_i\}$ is finite. Finite prime avoidance
gives an element $y$ of the prime outside every $\mathfrak m_i$ for
$i\in T$. Only finitely many variable blocks occur in $x,y$.
For each of these blocks at least one of $x,y$ is nonzero modulo its
$\mathfrak m_i$, so at most one scalar $c\in k$ has
$x+cy\in\mathfrak m_i$. Choose $c$ outside this finite set and also
so that $x+cy\ne0$. For a block not occurring in the two polynomials,
the same nonzero polynomial remains after setting that block to zero.
Thus $x+cy\in S$, a contradiction. The zero prime causes no problem.
The ideals $\mathfrak m_i$ are pairwise incomparable, so this shows
that the maximal ideals of $A$ are exactly $\mathfrak m_iA$.

Writing $K_i=k(t_j:j\ne i)$, we have
\[
 A_{\mathfrak m_iA}\simeq
 K_i[t_i^3,t_i^4,t_i^5]_{(t_i^3,t_i^4,t_i^5)}.
\]
Indeed localizing inverts every nonzero polynomial in the other blocks
and then every element not in the displayed maximal ideal. This is a
one-dimensional noetherian local domain, hence CM.
To prove $A$ itself noetherian, let $I\ne0$ be an ideal and choose
$0\ne x\in I$. Only finitely many maximal ideals contain $x$.
At each of them choose finitely many local generators of $I$ and lift
their numerators into $I$. Let $J$ be the ideal generated by all these
elements and by $x$. At the selected maximal ideals $J$ and $I$ agree;
at every other maximal ideal both localizations are the whole ring.
Thus $(I/J)_\mathfrak m=0$ for every maximal ideal, so $I=J$.
Every ideal is finite. Also $\dim A=1$ and all local rings are CM.
Every integral closed subscheme is either $\operatorname{Spec}A$ or
a closed point. Hence CMU holds, in fact in an especially strong form.

Nevertheless every nonempty principal open contains all but finitely
many of the maximal ideals $\mathfrak m_iA$, and at each such point
the parameter quotient by $t_i^3$ has basis
$1,t_i^4,t_i^5$. Its maximal ideal is square zero and its socle has
dimension two over $K_i$. Its injective hull has one-dimensional
socle, so the quotient ring is not that hull.
If a finite module $\Omega$ had Corollaire 10(b) on a dense open,
its generic rank would be one, by applying (b) at the generic point.
A finite rank-one module over a domain becomes free of rank one on
some nonempty principal open: choose a generic basis and kill the
finite kernel and cokernel of the resulting map. Choose a maximal
ideal in the intersection with the stipulated dense open. There
$\Omega/(t_i^3)\Omega$ is the quotient ring just computed, contradicting
the one-dimensional-socle condition. All hypotheses of the CMU extension
are therefore satisfied and its conclusion fails.

\textbf{Corrected statement.} Replace CMU by Gor-2, or retain either
of the two valid hypotheses ``locally embeddable in a regular scheme''
and ``excellent''. H035.19--20 prove these versions. Gor-2 is also
necessary for the corresponding assertion on every integral closed
subscheme: on a domain a generic canonical module is generically free
of rank one; wherever it is canonical and free the local ring is
Gorenstein. Thus the distinction at issue is sharp for the uniform
all-quotients use required by Théorème 12.

\section{H035.22: the absolute-value definition on the inserted typescript}
\textbf{Placement.} This separate arithmetic leaf is not about H035;
it is included solely because the six supplied entries discuss it.

\textbf{Literal defect and correction.} The written axioms say that
$f:K\to\mathbb R_+$ is multiplicative, subadditive, and
$f(x)=0\Rightarrow x=0$. They do not state the converse. The constant
function $f=1$ satisfies all those literal conditions, but does not
give the intended distance $d(x,y)=f(x-y)$, because $d(x,x)=1$.
Add $f(0)=0$. For $x\ne0$, positivity and multiplicativity now give
$f(1)=1$, $f(x^{-1})=f(x)^{-1}$ and $f(-1)=1$. Thus $d$ is positive
off the diagonal, symmetric, and satisfies the triangle inequality;
addition and multiplication are continuous by the triangle inequality
and $f(xy)=f(x)f(y)$. Inversion is continuous off zero because
$f(x^{-1}-y^{-1})=f(x-y)/(f(x)f(y))$ and $f(y)$ stays bounded
away from zero when $y$ is close to a fixed nonzero $x$.
For an additive real-valued valuation $v$ with $v(0)=+\infty$,
$f(x)=e^{-v(x)}$ has these properties and the stronger maximum
inequality, by the valuation axioms. This proves the construction as
well as the missing normalization.

\section{H035.23: the valuative criterion and binomial argument}
\textbf{Our normalized theorem.} For a normalized absolute value on
a field the following are equivalent: it comes from a real additive
valuation; $f(x+y)\leq\max(f(x),f(y))$; and $f(n\cdot1)\leq1$
for every positive integer $n$.

\textbf{Proof.} The valuation formulation and the maximum inequality
are interchanged by $v=-\log f$ on $K^\times$ and $v(0)=+\infty$.
The maximum inequality implies the integer bound by repeated addition
of $1$. Conversely the integer bound applies to every binomial
coefficient. Writing $C=\max(f(x),f(y))$, multiplicativity and the
ordinary triangle inequality give, for every positive $n$,
\[
 f(x+y)^n=f((x+y)^n)
 \leq\sum_{i=0}^n f\bigl({n\choose i}\bigr)f(x)^if(y)^{n-i}
 \leq(n+1)C^n.
\]
If $C=0$ both elements vanish. Otherwise take $n$th roots and let
$n\to\infty$, obtaining $f(x+y)\leq C$. If ``bounded on the
integers'' replaces ``at most one'', apply boundedness to $n^r$:
$f(n)^r\leq C_0$ for every $r$, hence $f(n)\leq1$. This verifies
the stronger boundedness wording used in the note entry.

\section{H035.24: positive characteristic and the root-of-unity slip}
\textbf{Our counterexample to the literal supporting sentence.}
In characteristic $p$, take $n=p$. Then $n\cdot1=0$, and zero is
not a root of unity. The printed claim about \emph{every} $n\cdot1$
is therefore false.

\textbf{Corrected result and proof.} Every nonzero element of the
prime field $\mathbb F_p$ has finite multiplicative order. Therefore
$f(a)^r=f(1)=1$ for some $r>0$, and $f(a)=1$. Zero has value zero
by the corrected definition. All integer multiples of $1$ have value
at most one, so H035.23 proves that every normalized absolute value
in positive characteristic is valuative. The mathematical corollary
survives the missing word in its supporting sentence.

\section{H035.25: equivalence of absolute values}
\textbf{Literal defect.} A real exponent without a positivity
qualification does not define the intended equivalence relation.
Exponent zero on $K^\times$ takes any absolute value to the trivial
one; this relation is not symmetric for a nontrivial original value.
At zero it moreover invokes the undefined expression $0^0$.
Negative exponents are undefined at zero and do not generally preserve
the triangle inequality. The correct exponent is strictly positive,
with both functions assumed to be absolute values.

\textbf{Our proof of the corrected equivalence and topology criterion.}
Positive-power equality is reflexive, symmetric and transitive by
exponents $1$, $1/s$ and $st$. It preserves the topology, since
$f(x)^s<\varepsilon$ iff $f(x)<\varepsilon^{1/s}$.
Conversely suppose two absolute values $f,g$ induce the same topology.
If one is trivial, its topology is discrete. A nontrivial absolute
value has an element $a$ with $0<f(a)<1$ (replace an element by its
inverse if necessary), and $a^n\to0$, so its topology is not discrete.
Hence both are trivial in that case. Otherwise convergence of $a^n$
shows that $f(a)<1$ iff $g(a)<1$. In particular $f(a)=1$ iff
$g(a)=1$, using the corresponding assertion for the inverse.
Fix $b$ with $f(b)>1$, so $g(b)>1$. For arbitrary $a\ne0$, apply
the equivalence $f(c)<1\Leftrightarrow g(c)<1$ to $c=a^n b^{-m}$,
where $n>0$ and $m\in\mathbb Z$. It says that the rational cuts
of $\log f(a)/\log f(b)$ and $\log g(a)/\log g(b)$ coincide.
They are equal, so $g(a)=f(a)^s$ for
$s=\log g(b)/\log f(b)>0$. Equality at zero follows from
normalization. This proves the full usual equivalence criterion.

\section{H035.26: the classification on the rational field}
\textbf{Corrected statement.} A normalized absolute value on
$\mathbb Q$ is either trivial, $p$-adic to a positive power for a
unique prime $p$, or the ordinary absolute value to a power
$0<s\leq1$. The trivial case must be included unless nontriviality
is stated. The typescript ends in the middle of the archimedean
argument; the completion below is ours.

\textbf{Nonarchimedean proof.} By H035.23 all integers have value at
most one. If all primes have value one, multiplicativity makes the
absolute value trivial. Otherwise choose $p$ with $f(p)<1$. No
distinct prime $q$ has $f(q)<1$: Bezout gives $ap+bq=1$, while the
maximum inequality would give
$1\leq\max(f(a)f(p),f(b)f(q))<1$. Thus all other primes have
value one. Factoring numerator and denominator yields
$f(x)=f(p)^{v_p(x)}=|x|_p^s$, with
$s=-\log f(p)/\log p>0$. Conversely these are absolute values,
and their small-valued prime identifies $p$ uniquely.

\textbf{Archimedean proof.} There is an integer $b\geq2$ with
$f(b)>1$. For an integer $a\geq2$, expand $a^n$ in base $b$:
$a^n=\sum_{i=0}^r c_i b^i$, where $0\leq c_i<b$ and
$r\leq n\log a/\log b$. With $C_b=\max_{0\leq c<b}f(c)$,
\[
 f(a)^n\leq C_b\sum_{i=0}^r f(b)^i
 \leq \frac{C_b f(b)}{f(b)-1}
              f(b)^{n\log a/\log b}.
\]
Taking $n$th roots gives $f(a)\leq a^s$, where
$s=\log f(b)/\log b>0$. Every integer $a\geq2$ actually has
$f(a)>1$. Otherwise expanding $b^n$ in base $a$ bounds its value
by $C_a(r+1)$, which grows only linearly in $n$, contradicting the
exponential $f(b)^n$. Interchanging $a,b$ in the preceding estimate
therefore gives the reverse logarithmic inequality and $f(a)=a^s$.
Multiplicativity extends this to all nonzero rationals, with symmetry
for negatives. The triangle bound $f(n)\leq n$ forces $s\leq1$.
Conversely $(u+v)^s\leq u^s+v^s$ for nonnegative $u,v$ and
$0<s\leq1$ proves the triangle inequality for $|\cdot|^s$:
divide by $(u+v)^s$ and use $t^s\geq t$ on $[0,1]$.
This proves both directions of the classification directly. The separate
equivalence-of-topologies statement has also been checked in Bourbaki,
\emph{Topologie générale}, IX, \S3, no.2, Proposition 6,
printed pp.TG IX.29--30; that proposition is not the classification
of absolute values on $\mathbb Q$.

\section{H035.27: detection by all parameter quotients (Corollaire 10c)}
\textbf{Our proof and strengthening.} Let $A$ be CM local, let
$\Omega$ have the parameter-injective property, and let $u:R\to R'$
be a map of finite modules. Suppose every parameter ideal $q$ makes
\[
 \operatorname{Hom}(R',\Omega/q\Omega)\longrightarrow
 \operatorname{Hom}(R,\Omega/q\Omega)
\]
an isomorphism. Artinian Matlis duality identifies this condition with
$R/qR\xrightarrow\sim R'/qR'$. The cokernel of $u$ is finite and
vanishes modulo one $q\subset\mathfrak m$, so Nakayama first gives
surjectivity. Its kernel $K$ is then contained in $qR$ for every
parameter ideal. Fix parameters $t_1,\ldots,t_d$ and use
$q_n=(t_1^n,\ldots,t_d^n)\subset\mathfrak m^n$. Krull intersection
gives $K\subset\bigcap_n\mathfrak m^nR=0$. Thus $u$ is an
isomorphism; the manuscript's assumed surjectivity can be omitted.
In dimension zero the empty parameter system has $q=0$ and Matlis
duality already concludes. A single parameter ideal is insufficient:
over the DVR $A=k[[t]]$, $A\twoheadrightarrow A/(t)$ is an
isomorphism modulo $(t)$ but is not an isomorphism. This also explains
why the universal parameter quantifier is retained.

\section{H035.28: the cancelled Corollaire 11 on scan 13}
\textbf{Reconstruction of the cancelled draft.} Read the final local
definitions as $\Phi_R(M)=\operatorname{Hom}(R,M)$ and similarly
for $R'$, retaining the cancellation of the whole draft. Assume these
two functors commute, on $\Omega$, with all parameter quotients, and
that $\operatorname{Hom}(u,\Omega)$ is an isomorphism, where
$u:R\twoheadrightarrow R'$ is a surjection of finite modules.

\textbf{Proof.} Tensor the last isomorphism with each $A/q$ and use
the two base-change isomorphisms. The resulting isomorphism is exactly
$\operatorname{Hom}(u,\Omega/q\Omega)$. H035.27 gives $u$ an
isomorphism. If one instead retains the cancelled formulas with an
additional $\Omega\otimes M$, this is a different assertion; the
proof here is for the final local rewrite actually shown in the scan.
Cancellation is an editorial fact, not evidence that the corrected
mathematical implication is false.

\section{H035.29: stationarity and the lost strictness hypothesis}
\textbf{Literal counterexample.} In the rewritten statement on scan
15 the words \emph{filtrant strict} and coherent terms occur only in
the cancelled preceding draft. Without them take $A=k$, $\Omega=k$,
$P_i=k$ for $i\geq1$, and every transition $P_j\to P_i$ for $j>i$
zero. The prorepresented functor is
$F(M)=\operatorname{colim}(M\xrightarrow0 M\xrightarrow0\cdots)=0$.
Thus $F(\Omega)=0$ is finite, and the only parameter quotient of the
field is itself, so condition (ii) holds. The displayed conclusion
$P_i\to P_{i_0}$ is an isomorphism for all $i\geq i_0$ fails.
The system is pro-isomorphic to zero, but the manuscript explicitly
asserts eventual isomorphisms of its transition maps; these assertions
are different. Dropping finite terms also invalidates the claim that
finite generation is equivalent to stationarity of the Hom system.

\textbf{Corrected theorem and proof.} Let $(P_i)$ be a cofiltered
strict system of finite $A$-modules, with surjective transitions, over
the affine $U$ of Corollaire 10. Suppose $F(\Omega)$ is finite and
the displayed parameter comparisons for $F$ are isomorphisms. The
modules $\operatorname{Hom}(P_i,\Omega)$ form an increasing union
inside $F(\Omega)$. Finitely many generators lie in one stage $i_0$,
so that stage and every subsequent stage equal $F(\Omega)$.
For any $p,q$, consider the square whose upper row is
\[
 F(\Omega)\otimes A_p/q\ \xrightarrow\sim\ F(\Omega_p/q\Omega_p)
\]
and lower row is the same natural map for $\operatorname{Hom}(P_{i_0},-)$.
The left vertical arrow is an isomorphism; the right vertical arrow
$\alpha$ is injective by strictness. Commutativity makes $\alpha$
surjective, since the composite of the lower horizontal and right
vertical arrows is the upper horizontal after the left isomorphism.
Thus $\alpha$ is an isomorphism. Every subsequent stage injects into
the same upper-right module, hence
$\operatorname{Hom}(P_{i_0},\Omega_p/q\Omega_p)\to
\operatorname{Hom}(P_i,\Omega_p/q\Omega_p)$ is an isomorphism.
H035.27 makes $(P_i)_p\to(P_{i_0})_p$ an isomorphism for every $p$.
A map of modules that is an isomorphism at all primes is an isomorphism.
This proves eventual stationarity. Notice that no unproved base-change
property for the individual $P_i$ was used; it follows from the square.

\section{H035.30: Théorème 12, the central representability criterion}
\textbf{Precise statement in his order.} Let $A$ be a noetherian
quotient of a regular noetherian ring, and $F:\operatorname{Mod}A\to
\mathrm{Sets}$. The reordered numbering on the page is:
\begin{enumerate}
\item[(i)] $F$ preserves filtered colimits.
\item[(ii)] $F$ is left exact (preserves finite limits).
\item[(iii a)] For every prime $p$ and finite $\widehat{A_p}$-module
$M$, the map $F(M)\to\lim_nF(M/\mathfrak m_p^{n+1}M)$ is bijective.
\item[(iii b)] For every $a\in A$, with $\widehat A$ the $(a)$-adic
completion, and every finite $\widehat A$-module $M$, the analogous
map $F(M)\to\lim_nF(M/a^{n+1}M)$ is bijective.
\item[(iv)] $F(M)$ is finite over $A$ for every finite $A$-module $M$.
\item[(v)] $p\mapsto\dim_{k(p)}F(k(p))$ is constructible.
\item[(vi)] For every nonzero quotient $B$ of $A$ and every finite
$B$-module $\Omega$, there is a nonnilpotent $a\in B$ such that
for all $p\in D(a)$ and every system of parameters $\mathbf x$ of
$B_p$, the canonical map
\[
 F(\Omega)\otimes_B B_p/(\mathbf x)
 \longrightarrow F\bigl(\Omega\otimes_B B_p/(\mathbf x)\bigr)
\]
is bijective.
\end{enumerate}
Then $F$ is represented by a finite $A$-module if and only if all
these conditions hold. The module structures and the canonical maps
come from (ii), as proved in H035.2. ``Nonzero'' in (vi) is an actual
insertion in the scan, not our repair.

\textbf{Our necessity proof.} For $F=\operatorname{Hom}(P,-)$ with
$P$ finite, (i), (ii) and (iv) are immediate from finite presentation
over a noetherian ring. Hom commutes with all inverse limits, and its
restriction to a completed ring is represented by the scalar extension
of $P$, proving both parts of (iii). A finite presentation matrix
shows that $\dim\operatorname{Hom}(P,k(p))$ is constructible, proving
(v). Over a quotient $B$, choose a nonempty principal CM open using
H035.19. Corollaire 7, H035.13, applied to $P\otimes B$ and the
given $\Omega$ gives a common dense open for every local regular
sequence. Parameter systems on the CM open are regular, which gives
(vi) after a further nonempty principal shrinking.

\textbf{Our sufficiency proof.} All six conditions pass to a quotient
$B=A/I$. For finite values use that finite $B$-modules are finite over
$A$; for residue ranks restrict the constructible function; for (vi)
use quotients of quotients. The completion assertions follow from
exactness of noetherian completion:
$\widehat B=\widehat A/I\widehat A$, including the localized version.
By H035.5 it suffices, for each nonzero quotient, to find a
nonnilpotent principal localization where its functor is represented.
Replace $A$ by that quotient.

Conditions (i), (ii), (iii a), (iv) imply the local representability
of H035.7. The finite-value condition (iv) supplies the hypothesis
missing from Remarque 5 bis. Therefore H035.10 applied to every
integral quotient, together with (v), gives generic representability
on each such quotient. H035.8 now supplies a strict prorepresentation
$F=\operatorname{colim}_i\operatorname{Hom}(P_i,-)$ with finite terms.
Use Corollaire 10 on all of $\operatorname{Spec}A$: it gives a finite
$A$-module $\Omega$ canonical on a \emph{dense} open $U$, meeting
every component. Condition (iv) makes $F(\Omega)$ finite. Apply (vi)
to this module and $B=A$, obtaining a nonempty $D(a)$. Density makes
$U\cap D(a)$ nonempty; choose a nonempty principal open $D(b)$
inside it. The localized strict system and $\Omega_b$ satisfy
H035.29: $F(\Omega_b)=F(\Omega)_b$ by (i), and the parameter
comparisons are exactly those of (vi) at the points of $D(b)$.
No claim about commuting localization with completion is needed here.
On this affine open H035.29 applies and makes the strict system
stationary. Its constant finite module represents the localized
functor. This is exactly the generic representation needed by H035.5,
which concludes over $A$.

The proof does not use the false CMU marginal enlargement, the false
unqualified converse in 5 bis, the erroneous $Z^1$ Tor criterion, or
the non-strict reading of Corollaire 11. Each dependency used here has
been proved with the hypotheses supplied by this theorem.

\section{H035.31: Remarque 13 and redundancy of the rank condition}
\textbf{Our statement.} Suppose instead that $A$ is noetherian and
every integral quotient $B$ has a nonempty regular principal open.
Then the criterion H035.30 holds without the regular-quotient
hypothesis, and (v) follows from (i), (ii), (iv), (vi).

\textbf{Proof.} The ring hypothesis implies Gor-2, so every use of
Corollaire 10 in H035.30 remains valid by H035.19. To obtain (v),
work on an integral quotient $B$, take $\Omega=B$ in (vi), and
intersect the resulting open with a regular open. At each point $p$
of this intersection a regular system of parameters generates the
maximal ideal of $B_p$. Thus (vi) identifies
$F(B)\otimes_B k(p)$ with $F(k(p))$. The module $F(B)$ is finite
by (iv), and becomes free on a smaller nonempty open. Its fibre
dimension is constant there. Noetherian induction, as in H035.9,
gives constructibility on all of $\operatorname{Spec}A$. The proof
of H035.30 now applies unchanged. This proves exactly the redundancy
claimed, without attributing the extra explanatory argument to the page.

\section{H035.32: proper Hom representability (Théorème 14)}
\textbf{Precise statement.} Let $f:X\to S$ be proper and finitely
presented, and let $F,G$ be finitely presented quasi-coherent modules
on $X$, with $G$ flat over $S$. There is a finitely presented module
$P$ on $S$, unique up to the unique isomorphism respecting the
representation, such that for every quasi-coherent $M$ on $S$,
\[
 \mathcal Hom_S(P,M)\simeq
 f_*\mathcal Hom_X(F,G\otimes_{\mathcal O_S}M).
\]
Formation of this representing object commutes with arbitrary base
change. The notation on the right includes the pullback of $M$.

\textbf{Exact literature certificate and full reduction.} Hall,
\emph{Cohomology and base change for algebraic stacks}, final author
version dated 30 April 2014, Theorem C (p.4), makes
$M\mapsto\operatorname{Hom}_X(F,G\otimes M)$ a coherent functor
over an affine base. Its hypotheses are locally finite presentation of
$f$, finite presentation and proper support of $G$, and flatness of
$G$ over the base; all hold here. Schemes are included among its
algebraic stacks. The source object is allowed to be an arbitrary
quasi-coherent module. Example 3.10 (p.18) says that a left exact
coherent functor is represented by a module, finitely presented when
the functor preserves filtered colimits. The functor here is left
exact because $G$ is flat over the base. It preserves filtered
colimits because $F$ is finitely presented and $X$ is quasi-compact
and quasi-separated over an affine base. Thus it is $\operatorname{Hom}(P,-)$
with $P$ finitely presented. Theorem D (p.4) independently packages
the resulting Hom functor as an abelian cone, affine and finitely
presented in this case.

Here are the base-change and sheaf details. For affine $S=\operatorname{Spec}A$
and $T=\operatorname{Spec}B\to S$, the projection $g:X_T\to X$ is
affine. For a $B$-module $M$, affine adjunction gives
\[
 \operatorname{Hom}_{X_T}(F_T,G_T\otimes_B M)
 =\operatorname{Hom}_X(F,G\otimes_A M)
 =\operatorname{Hom}_B(P\otimes_A B,M).
\]
The first equality follows by pushing the coefficient along $g$;
its underlying $\mathcal O_X$-module is $G\otimes_A M$, even if
$B$ is not flat. Yoneda therefore identifies the representing module
over $T$ with $P\otimes_A B$. These identifications satisfy cocycles
by uniqueness, glue over affine covers, and prove arbitrary base
change. Checking sections on all affine opens gives the displayed
sheaf isomorphism. Uniqueness means uniqueness together with this
functorial isomorphism, not that $P$ has no automorphisms.

\textbf{Historical qualification.} The opened EGA III$_2$ 7.7.8,
printed p.70, has an additional two-term locally free presentation
hypothesis on the source sheaf. Remark 7.7.9(iii), p.71, announces
its later removal. Consequently that corollary alone is not cited as
the full theorem written here. Hall's result matches the full hypotheses.

\section{H035.33: a quasi-coherent first argument}
\textbf{Our proof.} Keep the hypotheses on $f,G$ in H035.32, but
allow $F$ to be any quasi-coherent module. Hall Theorem C still
applies, and left exactness still holds. Example 3.10 therefore gives
a representing module $P$, without a finite-presentation conclusion.
The affine-adjunction calculation in H035.32 did not use finiteness
of $F$ or $P$, so also proves arbitrary base change and gluing here.

The passage to the limit proposed in the margin gives the same object.
Locally on $S$, $X$ is quasi-compact and quasi-separated. Write
$F=\operatorname{colim}_i F_i$ with the $F_i$ finitely presented,
using the finite-presentation approximation theorem for quasi-coherent
modules on such schemes. If $P_i$ represents the corresponding
functor, functoriality and Yoneda give maps in the same direction
$P_i\to P_j$ as $F_i\to F_j$. Then
\[
 \operatorname{Hom}_X(F,G\otimes M)
 =\lim_i\operatorname{Hom}_X(F_i,G\otimes M)
 =\lim_i\operatorname{Hom}(P_i,M)
 =\operatorname{Hom}(\operatorname{colim}_i P_i,M).
\]
Thus $P=\operatorname{colim}_iP_i$. This also shows why finiteness
is not retained. For example $f$ the identity of a field, $G=k$,
and $F$ an infinite-dimensional vector space gives $P=F$.

\section{H035.34: the higher underived coefficient functors}
\textbf{Precise reconstruction.} Let $A$ be noetherian and a quotient
of a regular noetherian ring, $f:X\to\operatorname{Spec}A$ proper,
and $F,G$ coherent on $X$. Fix nonnegative $i,j$. The additive
$A$-linear functors
\[
 T(M)=H^j\bigl(X,\mathcal Ext_X^i(F,G\otimes_A M)\bigr),
 \qquad T'(M)=\operatorname{Ext}_X^j(F,G\otimes_A M)
\]
satisfy H035.30(i), (iii a,b), (iv), (v), (vi). No flatness of $G$
is assumed. Tensor is \emph{underived}; replacing it by a derived
tensor would be a change of statement.

\textbf{Our proof of colimits and finiteness.} On a finite affine
cover of the separated noetherian $X$, and its affine intersections,
resolve $F$ by finite free modules. Computing a fixed Ext degree
uses only finitely many terms. Tensor, kernels, cokernels and filtered
colimits commute in modules. The finite Čech complex computes
cohomology of quasi-coherent sheaves, so the same statement holds
for $T$ and for the bounded part of the local-to-global Ext
spectral sequence computing $T'$. For finite $M$ the coefficient
and all relevant sheaf Ext modules are coherent. Proper coherence
of cohomology then makes $T(M),T'(M)$ finite. This proves (i),(iv).

\textbf{Our proof of the completion assertions.} Let $C$ be either
of the noetherian completions appearing in (iii), and let $M$ be
finite over $C$. Flat base change from $A$ to $C$ lets us compute
over $X_C$, with the coherent coefficient $E=G_C\otimes_C M$;
flatness is used only for this change of base, not for $G$.
Let $I$ be the ideal of completion. A bounded portion of a local
finite-free resolution of $F_C$ gives a complex $K$ of coherent
modules, with
$K/I^nK=\operatorname{Hom}(L,E/I^nE)$.
Artin--Rees applied to its cycles and boundaries proves that
\[
 \{H^i(K)/I^nH^i(K)\}_n\longrightarrow
 \{H^i(K/I^nK)\}_n
\]
is a pro-isomorphism. Explicitly, the filtrations on each cycle or
boundary induced from its adjacent finite module are cofinal with
its own $I$-adic filtration; tensoring the two exact sequences
for cycles and boundaries introduces kernels and cokernels whose
transition maps vanish after increasing $n$ by the Artin--Rees
constant. A finite affine cover supplies a common constant. This
argument works for any ideal of a noetherian ring and any bounded
finite complex; it makes no flatness assumption on its terms.

Consequently the sheaf Ext systems are pro-isomorphic to the
ordinary infinitesimal quotients of their coherent limiting sheaf.
Apply proper cohomology to this pro-isomorphism and then formal
functions, Stacks Theorem 30.20.5 (02OC), to obtain
$T(M)=\lim_nT(M/I^nM)$. The left module is finite over the complete
ring $C$, so equals its completion. Here is a more explicit justification
for total Ext, avoiding an unjustified interchange of inverse limits
with a spectral sequence. On finite $C$-modules the functor
$Q(V)=\{V/I^nV\}_n$ to pro-modules is exact: the assertion for a
submodule is precisely Artin--Rees. The local-to-global Ext spectral
sequence has finite $C$-modules on every page, and its part contributing
to any fixed total degree has a finite filtration and stabilizes after
finitely many pages. Apply the exact functor $Q$ to this spectral
sequence. Reduction of the coefficient induces a map to the pro-system
of local-to-global spectral sequences for $E/I^nE$; all these groups
are killed by $I^n$, so the maps factor through $Q$.
On the $E_2$ page the maps are pro-isomorphisms: first use the local
Artin--Rees comparison of sheaf Ext above, then the pro-version of
proper formal functions for each coherent sheaf. The latter says
that $\{H^p(X,H)/I^n\}_n\to\{H^p(X,H/I^nH)\}_n$ is a
pro-isomorphism; it is the Artin--Rees/Mittag--Leffler assertion in
the proof of Stacks Lemma 30.20.4 and Theorem 30.20.5.
Exactness in the pro-category carries the isomorphism through kernels,
images, the finitely many relevant pages, and the finite abutment
filtration. Thus
$\{\operatorname{Ext}^j(F_C,E)/I^n\}_n$ and
$\{\operatorname{Ext}^j(F_C,E/I^nE)\}_n$ are pro-isomorphic.
Taking limits now gives $T'(M)=\lim_nT'(M/I^nM)$.
This proves both actual
completed-module quantifiers in (iii), including modules not obtained
by completing a finite $A$-module.

\textbf{Our proof of constructibility.} Fix any integral quotient
$B$ of $A$. Keep the finite local free resolutions of the original
$F$ and tensor their \emph{coefficient} with $B$; do not assume that
a resolution of $F$ stays exact after this base change. The resulting
Hom complexes still compute the functors restricted to $B$-modules,
by their original definition. Their finitely many relevant terms,
cycles and cohomology are coherent over the finite type scheme $X_B$.
Generic freeness over the domain $B$ makes these terms and
cohomology flat on a nonempty base open. H035.36 then gives
sheaf-Ext base change for arbitrary modules on that open. The
resulting coherent sheaves are flat over the base after further
shrinking. Their finite Čech complexes are complexes of flat
$B$-modules. Properness makes their cohomology finite; generic
freeness makes this cohomology flat as well. A second application
of H035.36 yields $T(M)=T(B)\otimes_B M$ there. For $T'$ the
same argument is applied to the finite region of the double
complex, or equivalently its finite spectral sequence filtration.
In particular the residue-field ranks are constant on a smaller
nonempty open of each integral closed subset. Noetherian induction
gives (v).

\textbf{Our proof of (vi).} For a nonzero quotient $B$ and finite
$B$-module $\Omega$, apply H035.35 on the original $X$ to the
coefficient $G\otimes_A\Omega$, regarded over $\operatorname{Spec}B$.
The same local resolutions of the original $F$ justify this restriction.
Apply H035.38 to the resulting coherent Ext sheaves and, for $T'$,
to their bounded hypercohomology calculation. A common dense open
gives base change for every $B_p$-regular sequence. Intersect with
a CM open, available from H035.19 for quotients of $A$. There
every parameter system is regular. A nonempty principal subopen
proves (vi). This completes all the non-left-exact axioms claimed.

These functors need not be left exact: for $X=S=\operatorname{Spec}A$,
$G=A$, take $T'(-)=\operatorname{Ext}^1_A(A/(t),-)$ over
$A=k[t]$. It is $M\mapsto M/tM$, which sends the injection
$A\xrightarrow{t}A$ to the zero map on the nonzero module $A/tA$.
Nor does the preceding proof assert Auslander coherence. Hall
Theorem E uses derived tensor, and Jaffe's 1994 arXiv version,
section 7, pp.28--29, gives a proper underived-cohomology
counterexample to module coherence. These distinctions are retained.

\section{H035.35: generic regular-quotient base change for Ext (Lemme 15)}
\textbf{Our proof.} Let $X$ be of finite type over noetherian $A$,
let $F,G$ be coherent, $\Omega$ finite over $A$, and fix $N$.
Choose a finite affine cover of $X$. On each affine choose a
resolution $L_\bullet$ of $F$ by finite free modules over its
coordinate ring and put $K=\operatorname{Hom}(L,G\otimes_A\Omega)$.
The same complex tensored with an $A$-module $M$ computes the
Ext sheaves with coefficient $G\otimes_A\Omega\otimes_A M$.
Discard terms in degrees $\geq N+2$. This does not change
cohomology through degree $N$ either before or after tensoring,
and leaves a bounded complex of coherent modules. Its terms and
cohomology form a finite list. H035.37 supplies a common dense
base open on which all positive Tor of these modules vanishes
against every quotient by a local $A$-regular sequence. H035.36
then gives the required comparison through degree $N$. A finite
intersection of the opens for the cover remains dense. Since the
comparison is canonical, its local isomorphisms give the asserted
sheaf isomorphism. If $N<0$ the assertion is vacuous.

\section{H035.36: the bounded-complex Tor criterion}
\textbf{Our proof.} Let $K^\bullet$ be a bounded cochain complex
over a ring $A$, and $M$ a module. Assume all positive Tor groups
of $M$ with every $K^j$ and every $H^j(K)$ vanish. Choose a free
resolution $L_\bullet\to M$. The double complex
$L_\bullet\otimes K^\bullet$ has only finitely many $K$-columns.
Taking homology first in the resolution direction leaves only
$M\otimes K^\bullet$, by the Tor assumption on the terms.
Taking it first in the $K$ direction is valid because each $L_i$
is free; the next page is
$\operatorname{Tor}_i^A(M,H^j(K))$. It is supported on $i=0$,
by the second Tor assumption. Both spectral sequences converge
because there are finitely many summands on each total diagonal.
Their identifications give
$H^j(K)\otimes M\simeq H^j(K\otimes M)$. The edge morphism
is the usual map induced by cycles, since the augmentation
$L_\bullet\to M$ is the map used in the first comparison.
Thus the result proves the canonical comparison, not just an
abstract equality of dimensions. This stronger criterion is valid
and should be distinguished from the false two-Tor criterion H035.15.

\section{H035.37: relative generic Tor vanishing (Lemme 16)}
\textbf{Reading.} The insertion names the coherent module $R$,
whereas the displayed Tor uses $F$. The coefficient inside the first
associated-graded expression is not securely distinguishable in the
crop: it may be a structure-sheaf symbol or a coefficient symbol.
We do not replace it on mathematical grounds. Our proof needs the
associated graded of the coefficient $R$ as well as that of $A$;
the definite $R/F$ mismatch and the unresolved graded symbol are
recorded separately in READINGS.tsv.

\textbf{Our proof.} Let $X$ be of finite type over noetherian $A$,
and let $R$ be coherent on $X$. Write $J=\sqrt{0_A}$. Remove
intersections of components of $\operatorname{Spec}(A/J)$.
Generic freeness on each remaining integral component makes
all finitely many pieces of $\operatorname{gr}_J A$ locally free,
and all pieces of $\operatorname{gr}_J R$ flat over $A/J$, on a
common dense base open. For the latter assertion use a finite
affine cover of $X$: the pieces are finite modules over finite
type $A/J$-algebras, so relative generic freeness applies. They
need not be finite as $A/J$-modules, and no such assertion is made.
At any base point of this open, the last nonzero free graded piece
of $A$ shows, as in H035.12, that an $A$-regular element remains
regular modulo $J$. Flatness makes it regular on each graded
piece of $R$. Strictness of the induced filtrations gives the same
statement after quotienting. Induction handles any regular
sequence. Its Koszul resolution then proves
$\operatorname{Tor}_i^A(R,A_p/(\mathbf x))=0$ for every $i>0$.
This verifies the coefficient as well as the structure-sheaf case.

\section{H035.38: the quantifier in Lemme 17}
\textbf{Literal counterexample.} The scan introduces a dense $U$
but then says ``$\forall s\in S$''. Let $A=k[t]$,
$X=\mathbb P^1_A$, and choose the extension
\[
 0\longrightarrow\mathcal O_X(-2)\longrightarrow E
 \longrightarrow\mathcal O_X\longrightarrow0
\]
with class $t$ in $\operatorname{Ext}^1(\mathcal O,\mathcal O(-2))
=H^1(\mathcal O(-2))=A$. Such an extension exists by the
Čech description on the two standard affine charts, and $E$ is
locally free, since an extension of locally free modules locally
splits. The familiar two-chart computation gives
$H^0(\mathcal O(-2))=0$, $H^1(\mathcal O(-2))=A$,
$H^0(\mathcal O)=A$, $H^1(\mathcal O)=0$. The connecting map
is multiplication by the chosen extension class $t$. Thus
$H^0(E)=0$ and $H^1(E)=A/(t)$. At $t=0$ the extension class
vanishes, so $E_0=\mathcal O\oplus\mathcal O(-2)$ and
$H^0(E_0)=k$. The regular parameter $t$ in $A_{(t)}$ therefore
gives $0=H^0(E)\otimes k\to H^0(E\otimes k)=k$, not an
isomorphism. The properness, noetherianity and coherence hypotheses
are all satisfied, even with $E$ locally free.

\textbf{Corrected theorem and proof.} Replace $s\in S$ by
$s\in U$. For proper $f:X\to\operatorname{Spec}A$, a finite
affine cover gives a finite bound on cohomology. The modules
$H^q(X,E)$ are finite. H035.12, applied to this finite list,
and H035.37, applied to $E$, give a common dense open where
all their positive Tor groups vanish against every local
regular-sequence quotient $M=A_p/(\mathbf x)$.
The Koszul complex resolving $M$ is finite free. Tensoring it
with $E$ is a resolution of the underived $E\otimes_A M$,
by the latter Tor vanishing. Apply $R\Gamma$ to this bounded
complex. Because its terms are finite direct sums of $E$, it
computes $R\Gamma(E)\otimes_A^{\mathbf L}M$. The resulting
spectral sequence has entries
$\operatorname{Tor}_i^A(H^q(X,E),M)$; only $i=0$ survives.
Its edge map is precisely
$H^q(X,E)\otimes M\to H^q(X,E\otimes M)$, so is an
isomorphism. The same finite-complex argument proves the
bounded coherent-complex version used in H035.34. This proof
does not require flatness of $E$ over the original base.

\section{H035.39: the constructible-functor definition and composition programme}
\textbf{His words.} ``foncteur (additif ?)''; ``composé de foncteurs
constructibles''.

\textbf{State of this item.} The page proposes, rather than defines,
a two-scheme class. For the affine interpretation we take additive
$A$-linear functors $\operatorname{Mod}A\to\operatorname{Mod}A$
satisfying H035.30(i), (iii), (iv), (v), (vi), dropping (ii).
This is an explicit reconstruction: after finite-limit preservation
is removed, a set-valued functor does not automatically acquire the
module structures needed to formulate the other axioms. The exact
composition assertion for this whole affine class is still being
checked. It is not declared proved merely because a smaller class
of coherent functors has the desired closure property.

\textbf{Current proved reduction.} Filtered colimits and finite values
are preserved by composition. For residue fields, additivity and
$k(p)$-linearity give
$\dim F(G(k(p)))=\dim F(k(p))\dim G(k(p))$, so constructible rank
is preserved. Condition (vi) can in fact be strengthened from a
nonempty open to a dense open, without assuming irreducibility.
Fix a nonzero quotient $B$ and a finite $B$-module $\Omega$.
For each minimal prime $p$ choose $f\notin p$ in every other
minimal prime; if $p$ is the only one, take $f=1$.
Let $I=\ker(B\to B_f)$ and $C=B/I$. Since $B$ is noetherian,
$I$ is killed by a power of $f$. The support of $C$ is $V(p)$:
its primes over $D(f)$ are dense there, and contraction from
$B_f$ gives exactly this closure. In particular $\operatorname{Spec}C$
is irreducible and $C_f=B_f$. Apply (vi) on $C$ to
$\Omega/I\Omega$. Its nonempty good open meets $D(f)$;
on that intersection it is the original comparison on $B$,
because localization commutes with either additive finitary functor.
Its inverse image inside $D(f)\subset\operatorname{Spec}B$
is an open neighbourhood of $p$. Taking the union over the finitely
many minimal primes gives a dense good open. Consequently the good
opens for $G$ at $\Omega$ and for $F$ at the finite module
$G(\Omega)$ have dense intersection. On this intersection the two
comparison isomorphisms compose to that for $F\circ G$.
This proves (vi) for composition, including reducible quotients.

The only remaining affine axiom is completion. The unresolved issue is
whether a finite module over a completion has finite image there
under $G$, or whether the outer functor is continuous on the
resulting pro-finite system without that finiteness. The conditions
as written assert finite values only for finite modules over $A$.
This section is deliberately a working state, not a proof certificate.

\section{H035.40: the proposed Ext, Tor and direct-image examples}
\textbf{Literal scope.} No hypotheses on a general morphism are
specified in the proposed two-scheme theory. If the proposed direct
image assertion is read without a finiteness hypothesis, it is false:
for $f:\mathbb A^1_k\to\operatorname{Spec}k$,
$f_*\mathcal O_{\mathbb A^1}=k[t]$ is not finite-dimensional.
It already fails the finite-value requirement in degree zero.

\textbf{Precise positive realization.} For finite modules over a
noetherian ring, $\operatorname{Ext}^i_A(P,-)$ and
$\operatorname{Tor}_i^A(Q,-)$, with fixed finite $P,Q$ and fixed
degree, are computed by a finite portion of a resolution by finite
free modules. They are additive, preserve filtered colimits, and
have finite values on finite modules. Finite-complex Artin--Rees
proves their completion comparisons exactly as in H035.34.
Generic freeness of the finitely many terms and cohomology, on
each integral quotient, proves constructible ranks. H035.12
and H035.36 prove the regular-quotient comparisons; on rings
with CMU the parameter version follows by a CM shrinking.
For the proper coherent-coefficient direct-image functors
$M\mapsto H^q(X,E\otimes_A M)$, H035.34 supplies all the
axioms, without assuming the coefficient flat. This proves the
natural proper version suggested by the immediately preceding
lemmas, while retaining the failure of the unqualified assertion.
The missing general two-scheme definition and its composition
claim are tracked in H035.39, not hidden in this positive version.

\section{H035.41: affine restriction of scalars (Théorème 19)}
\textbf{Precise statement.} Let $f:X\to S$ be proper, flat and
finitely presented, and $Z\to X$ affine. The functor on $S$-schemes
\[
 T\longmapsto\operatorname{Hom}_{X_T}(X_T,Z_T)
\]
is represented by an affine $S$-scheme, compatibly with arbitrary
base change. No finiteness assumption on the affine $Z\to X$
is needed for this assertion.

\textbf{Our construction and proof.} Write $Z=\operatorname{Spec}_X B$
for a quasi-coherent $\mathcal O_X$-algebra $B$. H035.33, with
$G=\mathcal O_X$ (flat over $S$), supplies a quasi-coherent module
$P_B$ such that $W=\operatorname{Spec}_S\operatorname{Sym}P_B$
represents all $\mathcal O_{X_T}$-linear maps $B_T\to\mathcal O_{X_T}$.
Over $W$ there is the universal such map $u$. It is an algebra map
precisely when the following two linear maps are zero:
\[
 u\circ 1_B-1_{\mathcal O_X}:\mathcal O_{X_W}\to\mathcal O_{X_W},
 \qquad
 u\circ\mu_B-\mu_{\mathcal O_X}\circ(u\otimes u):
 B_W\otimes B_W\to\mathcal O_{X_W}.
\]
H035.33 and its base-change property turn each into a homomorphism
from a quasi-coherent module on $W$ to $\mathcal O_W$. Let $J$
be the sum of their image ideals. The closed subscheme $V(J)\subset W$
has the required universal property: after any $T\to W$ the maps
vanish exactly when $J$ maps to zero in $\mathcal O_T$, equivalently
when the map factors through $V(J)$. This equivalence is valid for
nonflat $T$, since it concerns the images of the defining maps,
not an unjustified identification of a pulled-back ideal with an
injective submodule. Thus $V(J)$ represents unital multiplicative
linear maps, namely sections of $Z_T$. It is affine over $S$.
All constructions are universal for arbitrary $T$, so base change
and gluing on $S$ follow from Yoneda. Associativity imposes no
additional equation, because both source and target already are
associative algebras and multiplicativity is imposed on all pairs.

\section{H035.42: finite presentation and finite type in Théorème 19}
\textbf{Our direct proof.} Work over affine $S$; then $X$ is
quasi-compact and quasi-separated. Suppose first that $B$ is a
finite type $\mathcal O_X$-algebra. Write its underlying module
as a filtered colimit of finitely presented modules. A finite
affine cover and finitely many algebra generators show that one
stage $Q\to B$ has image generating $B$ as an algebra. Hence
$\operatorname{Sym}Q\twoheadrightarrow B$ and $Z$ is closed
in the linear cone $\mathbb V(Q)$. H035.32 makes the restriction
of this cone $\mathbb V(P_Q)$ with $P_Q$ finitely presented.
H035.43 makes the restriction of $Z$ a closed subscheme of it.
Since every closed immersion is of finite type, the restriction
of $Z$ is of finite type over $S$.

If $B$ is finitely presented as an algebra, the kernel
$I=\ker(\operatorname{Sym}Q\to B)$ is of finite type as an
ideal of $\operatorname{Sym}Q$. Indeed both algebras have finite
presentations over $\mathcal O_X$, so after introducing finite
generators the relative presentation has finitely many relations.
Again using a finite affine cover and finite-presentation
approximation of the underlying quasi-coherent module $I$,
choose a finitely presented $\mathcal O_X$-module $R\to I$
whose image generates $I$ as an ideal. On
$W=\mathbb V(P_Q)$, extend the universal linear map on $Q$
to an algebra map $\operatorname{Sym}Q\to\mathcal O_{X_W}$.
It factors through $B$ exactly when its composite with $R$
vanishes. H035.32 represents this composite by a homomorphism
$(P_R)_W\to\mathcal O_W$. Since $P_R$ is finitely presented,
its image ideal is finitely generated. Thus the required closed
subscheme is finitely presented over $W$, and hence over $S$.
This proves the stronger finiteness assertion directly, without
assuming the base noetherian or silently treating every closed
immersion as finitely presented.

The distinction is necessary. For the identity morphism of a
nonnoetherian affine scheme, restriction of scalars is the identity.
The closed subscheme defined by a non-finitely-generated ideal is
of finite type but is not finitely presented. For instance take
$A=k[x_1,x_2,\ldots]$ and $I=(x_1,x_2,\ldots)$.

\section{H035.43: preservation of closed immersions (Corollaire 20)}
\textbf{Our proof.} Let $Z\hookrightarrow Z'$ be a closed
immersion of affine $X$-schemes. By H035.41 both section functors
are represented. Pull back to the universal section of $Z'$ over
its representing scheme $W$. Its inverse image of $Z$ is a
closed subscheme $Y\subset X_W$. A base change $T\to W$ has
its universal section factor through $Z$ exactly when
$Y_T=X_T$. If $I\subset\mathcal O_{X_W}$ defines $Y$, use
H035.33 for $I$ to represent the inclusion by a map
$P_I\to\mathcal O_W$. The image ideal cuts out precisely the
bases where this map, and hence $I\to\mathcal O_{X_W}$,
vanishes. That condition is $Y_T=X_T$, including nonflat base
changes. It is therefore a closed subscheme of $W$, proving
the claimed closed immersion. This proof uses H035.41 but not
the finiteness assertion H035.42, so those dependencies are not
circular.

\section{H035.44: the linear-cone case and the missing pushforward}
\textbf{Our corrected formula and proof.} For a quasi-coherent
module $F$ on $X$, put $\mathbb V(F)=\operatorname{Spec}_X
\operatorname{Sym}F$, with this explicit convention. Its sections
over $X_T$ are $\operatorname{Hom}(F_T,\mathcal O_{X_T})$.
For $f$ proper flat finitely presented, H035.32--33 represent
these by $\mathbb V(P_F)$ on $S$. The sheaf formula is
\[
 \mathcal Hom_S(P_F,M)=
 f_*\mathcal Hom_X(F,\mathcal O_X\otimes_{\mathcal O_S}M).
\]
The displayed formula on scan 23 omits $f_*$, although its two
sides otherwise live on different schemes. The insertion above
is our correction, not a new reading of the ink.

When the manuscript calls $\mathbb V(F)$ a vector bundle for
coherent $F$, use the precise term \emph{linear cone}. It is a
vector bundle if $F$ is finite locally free. For example over
$A=k[t]$, $F=A/(t)$ gives $\operatorname{Spec}A[x]/(tx)$;
its fibre dimension jumps from zero to one and it is not a
vector bundle. The representability argument uses only the
linear cone, so this terminological correction changes no result.
Every affine finite type $X$-scheme embeds in such a cone with
$F$ finitely presented on a quasi-compact quasi-separated $X$,
by the generator argument in H035.42. In the noetherian setting
of the manuscript's reduction, ``coherent'' and ``finitely
presented'' coincide, proving exactly its embedding step.

\section{H035.45: the closed-subscheme condition and Murre's theorem}
\textbf{Exact historical reference.} The opened Murre exposition
294, \emph{Les schémas de Picard}, May 1965, Theorem 3, printed
p.255 (PDF page 14), treats the subfunctor of the base for a
closed $Y\subset X$, with $X$ flat of finite type over a locally
noetherian base. It is representable by a closed subscheme under
its adic condition, in particular when $X$ is proper. This is
the precise special case cited on scan 24, after its noetherian
reduction. It is not a claim that an arbitrary flat morphism has
this property.

\textbf{Our proof of the version needed here.} For proper flat
finitely presented $X\to S$, without noetherianity of $S$,
let $I$ be the quasi-coherent ideal of $Y$. H035.33 represents
$\operatorname{Hom}_X(I,\mathcal O_X\otimes M)$ by $P_I$.
The inclusion corresponds to $P_I\to\mathcal O_S$. Let $J$
be its image. For every $T\to S$, base change identifies
vanishing of this map with vanishing of the pulled-back ideal
map $I_T\to\mathcal O_{X_T}$, which says $Y_T=X_T$.
Thus $\operatorname{Spec}_S(\mathcal O_S/J)$ represents the
subfunctor, with either one element or no elements on each $T$.
This proves the claim and the required arbitrary base changes.
For comparison, SGA 3 VIII, Theorem 6.4 in the opened
Gille--Polo revision, pp.18--19, proves the same closed-condition
construction for an essentially free morphism. Its different
hypothesis is recorded rather than conflated with proper flatness.

\section{H035.46: restriction of a quasi-coherent linear cone (Remarque 21)}
\textbf{Our proof.} For proper flat finitely presented $f$, take
$G=\mathcal O_X$ in H035.33. For every quasi-coherent $F$ there
is a quasi-coherent $P$ on $S$ such that, for every $T\to S$,
\[
 \Gamma(X_T,\mathbb V(F)_T)
 =\operatorname{Hom}_{X_T}(F_T,\mathcal O_{X_T})
 =\operatorname{Hom}_T(P_T,\mathcal O_T)
 =\Gamma(T,\mathbb V(P)_T).
\]
Yoneda gives the required isomorphism of section functors.
If $F=\operatorname{colim}F_i$ with finitely presented terms,
then $P=\operatorname{colim}P_i$, and taking the relative
spectrum of the symmetric algebra changes these colimits into
the inverse limits of the cones. This is the passage to the
limit written on the page. There is no finite-type assertion
for arbitrary $F$; the identity-morphism example in H035.33
already shows why none should be added.
